% Reading list: the main papers cited in the project, with links and short summaries.
% Summaries are based on each paper's published abstract (ACL Anthology / arXiv).
\documentclass[11pt,a4paper]{article}
\usepackage[T1]{fontenc}
\usepackage[margin=2.2cm]{geometry}
\usepackage{times}
\usepackage{xcolor}
\usepackage{enumitem}
\usepackage[colorlinks=true,urlcolor=blue!60!black,linkcolor=black]{hyperref}
\usepackage{titlesec}
\titleformat{\section}{\large\bfseries}{}{0pt}{}
\titlespacing*{\section}{0pt}{1.2em}{0.4em}
\setlength{\parindent}{0pt}
\setlength{\parskip}{0.3em}

\definecolor{main}{HTML}{3A8F5C}
\definecolor{findings}{HTML}{E07B39}
\definecolor{preprint}{HTML}{7F8C8D}
\newcommand{\tagbox}[2]{\colorbox{#1!15}{\textcolor{#1!80!black}{\footnotesize\textbf{#2}}}}
\newcommand{\paper}[6]{% name, full title, authors/venue, tag, url, summary + relevance
  \subsection*{#1 \hfill #4}
  \textit{#2}\\
  {\small #3}\\
  {\small\url{#5}}\par
  #6
  \vspace{0.4em}}
\newcommand{\why}[1]{\par\textbf{Relevance to our project:} #1}

\title{Project 10: Code Reuse / Refactoring with AI\\[0.3em]
  \large Key papers: links and summaries}
\author{Evgenii Lialin \and Roman Vikhman}
\date{}

\begin{document}
\sloppy
\maketitle

\noindent Venue status under the course rule (A/A+/A* IR/NLP/CL venues, 2025 onward):
\tagbox{main}{main venue} counts;
\tagbox{findings}{Findings / Industry} needs the instructor's approval;
\tagbox{preprint}{preprint} arXiv, not peer-reviewed at a qualifying venue.
Summaries follow each paper's published abstract.

% ===========================================================================
\section{1.\quad The problem: models do not reuse existing code}

\paper{RepoExec}
{On the Impacts of Contexts on Repository-Level Code Generation}
{Hai et al.\ --- Findings of NAACL 2025}
{\tagbox{findings}{Findings}}
{https://aclanthology.org/2025.findings-naacl.82/}
{A benchmark for repository-level code generation that checks executability, functional correctness and
\emph{dependency utilisation}. Evaluating 18 models, it finds that full dependency context works best and
that small contexts can be misleading. Pretrained models are often correct but \emph{re-implement}
dependencies; instruction-tuned models use dependencies more but sometimes add needless complexity.
It introduces the \textbf{Dependency Invocation Rate (DIR)}, a metric of how often the generated code calls
the available dependencies, and an instruction-tuning dataset that raises DIR by over 10\%.
\why{The closest prior work. We adopt its central idea, measuring whether the model \emph{calls} the
existing function, and move it from generation to refactoring, where the duplicated logic is already in
the code and has to be recognised.}}

\paper{More Code, Less Reuse}
{More Code, Less Reuse: Investigating Code Quality and Reviewer Sentiment towards AI-generated Pull Requests}
{Huang et al.\ --- arXiv:2601.21276, 2026}
{\tagbox{preprint}{preprint}}
{https://arxiv.org/abs/2601.21276}
{An empirical study of pull requests written by LLM agents versus humans, looking beyond pass rates at
code-quality metrics and reviewers' reactions. LLM agents frequently \textbf{disregard code reuse
opportunities}, producing more redundancy than human developers, yet reviewers respond to AI pull requests
neutrally or even more positively. The authors argue that plausible-looking AI code hides redundancy,
so technical debt accumulates silently.
\why{Our main motivation: the lack of reuse is observed in real development and goes unnoticed in review.
Our benchmark asks where in the pipeline that reuse is lost.}}

\paper{DependEval}
{DependEval: Benchmarking LLMs for Repository Dependency Understanding}
{Du et al.\ --- Findings of ACL 2025}
{\tagbox{findings}{Findings}}
{https://aclanthology.org/2025.findings-acl.373/}
{A hierarchical benchmark built from 2,683 real repositories in 8 programming languages, with three
tasks: dependency recognition, repository construction and multi-file editing. Over 25 LLMs show
substantial gaps in repository-level understanding.
\why{Shows that understanding repository dependencies is a known weakness of LLMs; reuse requires exactly
this understanding.}}

% ===========================================================================
\section{2.\quad Refactoring benchmarks}

\paper{RefactorBench}
{RefactorBench: Evaluating Stateful Reasoning in Language Agents Through Code}
{Gautam et al.\ --- ICLR 2025}
{\tagbox{main}{ICLR 2025}}
{https://arxiv.org/abs/2503.07832}
{100 handcrafted multi-file refactoring tasks in popular open-source (Python) repositories, each given
by three natural-language instructions of different specificity. Agents solve only 22\% with the base
instructions, against 87\% for a human developer under time pressure. Trajectory analysis finds that
agents lose track of past actions; conditioning the agent on a representation of its state improves
results by 43.9\%.
\why{Python refactoring benchmark; it measures whether the task is completed, not whether existing code
is reused.}}

\paper{CodeTaste}
{CodeTaste: Can LLMs Generate Human-Level Code Refactorings?}
{Thillen et al.\ --- ICML 2026}
{\tagbox{main}{ICML 2026}}
{https://arxiv.org/abs/2603.04177}
{Starts from the observation that coding agents accumulate complexity, \textbf{duplication} and
architectural debt. A benchmark mined from large multi-file open-source refactorings, scored with test
suites plus static checks (dataflow-based) that verify removal of undesired and introduction of desired
code patterns. Agents implement refactorings well when they are specified in detail but often fail to
\emph{discover} the refactoring a human chose; a propose-then-implement decomposition helps.
\why{Closest in spirit: the gap is between executing a refactoring and choosing it. Our finding that
models ignore a helper already in context is a specific case of this ``choice'' failure.}}

\paper{SWE-Refactor}
{SWE-Refactor: A Repository-Level Benchmark for Real-World LLM-Based Code Refactoring}
{Xu et al.\ --- arXiv:2602.03712, 2026}
{\tagbox{preprint}{preprint}}
{https://arxiv.org/abs/2602.03712}
{1,099 developer-written, behaviour-preserving refactorings mined from 18 Java projects (922 atomic,
177 compound), each validated by compilation, tests and refactoring-detection tools. Nine LLMs are
evaluated; compound refactorings are the main source of failure (a Codex agent reaches 39.4\%).
\why{Our task validation follows the same principle (behaviour preserved, checked by tests), but we
target Python and reuse specifically.}}

\paper{MANTRA}
{MANTRA: Enhancing Automated Method-Level Refactoring with Contextual RAG and Multi-Agent LLM Collaboration}
{Xu et al.\ --- arXiv:2503.14340, 2025}
{\tagbox{preprint}{preprint}}
{https://arxiv.org/abs/2503.14340}
{An agent framework for method-level refactoring that combines context-aware retrieval-augmented
generation, multi-agent collaboration and verbal reinforcement learning. On 703 pure refactorings from
10 Java projects it produces code that compiles and passes tests in 82.8\% of cases (8.7\% for a plain
LLM baseline); developers judge its refactorings as readable and reusable as human ones.
\why{Shows that retrieval helps refactoring; we test whether retrieval is in fact the bottleneck for
reuse (it mostly is not).}}

% ===========================================================================
\section{3.\quad Retrieval of repository code}

\paper{CodeRAG}
{CodeRAG: Finding Relevant and Necessary Knowledge for Retrieval-Augmented Repository-Level Code Completion}
{Zhang et al.\ --- EMNLP 2025}
{\tagbox{main}{EMNLP 2025}}
{https://aclanthology.org/2025.emnlp-main.1187/}
{Addresses poor query construction, single-path retrieval and the mismatch between retriever and code LLM
in repository-level completion, using log-probability-guided queries, multi-path retrieval and
preference-aligned reranking (BestFit). Outperforms prior methods on ReccEval and CCEval.
\why{Representative of work that improves \emph{retrieval}. Our funnel shows that, for reuse, the
model's decision costs more than retrieval misses.}}

\paper{CoRet}
{CoRet: Improved Retriever for Code Editing}
{Fehr et al.\ --- ACL 2025 (short)}
{\tagbox{main}{ACL 2025}}
{https://aclanthology.org/2025.acl-short.62/}
{A dense retriever for code editing that combines code semantics, repository structure and call-graph
dependencies, trained with a loss designed for repository-level retrieval. Improves retrieval recall by
at least 15 points on SWE-bench and Long Code Arena bug localisation.
\why{Retrieval for \emph{editing}, our setting; a candidate for the dense-retrieval condition we have not
run yet.}}

\paper{ReflectCode}
{Beyond Function-Level Search: Repository-Aware Dual-Encoder Code Retrieval with Adversarial Verification}
{Liu et al.\ --- Findings of EMNLP 2025}
{\tagbox{findings}{Findings}}
{https://aclanthology.org/2025.findings-emnlp.1147/}
{Introduces RepoAlignBench for repository-level retrieval driven by change requests, and ReflectCode,
a dual-encoder (code and documentation towers) that uses an LLM to integrate syntax, function
dependencies and the intent of the change. Improves top-5 accuracy by 12.2\% and recall by 7.1\%.
\why{Retrieval driven by a change request, close to our refactoring request.}}

\paper{Think-Search-Patch}
{Think-Search-Patch: A Retrieval-Augmented Reasoning Framework for Repository-Level Code Repair}
{Xiong et al.\ --- EMNLP 2025 Industry Track}
{\tagbox{findings}{Industry}}
{https://aclanthology.org/2025.emnlp-industry.109/}
{A three-stage framework for repository-level repair: break the task down and write a search query,
retrieve snippets (e.g.\ with E5), then generate a patch; trained with SFT and DPO. Improves retrieval
accuracy and repair success on SWE-bench Lite, including multi-file bugs.
\why{Another retrieve-then-edit pipeline, which assumes the model uses what is retrieved; our results
question that assumption for reuse.}}

% ===========================================================================
\section{4.\quad How models use context and instructions}

\paper{Lost in the Middle}
{Lost in the Middle: How Language Models Use Long Contexts}
{Liu et al.\ --- TACL 2024}
{\tagbox{main}{TACL 2024}}
{https://aclanthology.org/2024.tacl-1.9/}
{On multi-document question answering and key-value retrieval, performance depends strongly on the
position of the relevant information: it is highest at the beginning or end of the context and drops
in the middle, even for long-context models.
\why{Basis of our position ablation: we place the helper first or last among the retrieved snippets.
Three of four models change their reuse rate with the position.}}

\paper{CodeIF}
{CodeIF: Benchmarking the Instruction-Following Capabilities of Large Language Models for Code Generation}
{Yan et al.\ --- ACL 2025 Industry Track}
{\tagbox{findings}{Industry}}
{https://aclanthology.org/2025.acl-industry.89/}
{The first benchmark of how well LLMs follow task-oriented instructions in code generation, covering
function synthesis, debugging, algorithmic refactoring and code explanation across difficulty levels and
domains; it highlights gaps in producing consistent, maintainable code.
\why{Motivates our two instructions: a generic refactoring instruction versus an explicit ``reuse
existing functions'' instruction, which raises reuse in every condition.}}

\end{document}
